%% sections/05-discussion.tex

\section{Discussion}
\label{sec:discussion}

\subsection{Why the Spine Works: Structural Analysis}

The campaign spine's 100\% completion rate and 71\% Route D/E rate cannot be
attributed to any single component. The five components function as a system:
the central question gives the hints their convergence target; the hints give
the routes their achievability conditions; the spotlight schedule gives the
deliverable its raw material; and the deliverable gives the finale its
non-scripted emotional content. Remove one component and the system degrades;
the corpus shows no campaigns designed with partial spine implementations,
but the component-interaction analysis in Section~\ref{sec:components}
predicts specific failure modes for each absent component.

The most critical component for reliable \textit{completion} is the branching
end-conditions. Campaigns without explicit completion criteria cannot
self-terminate, and in AI-simulated play this produces indefinite continuation
rather than a clean ending. The route structure guarantees that the campaign
can end gracefully at any engagement level; Routes A--C provide the structural
floor.

The most critical component for \textit{quality} of completion --- as measured
by Route D/E achievement and deliverable richness --- is the four-hint schedule.
Without Hint 3 and Hint 4, the party has no basis for recognizing that the
optimal ending is available. The hints do not force the party to reach Route D/E;
they create the structural possibility for it. A party that reaches the finale
with all four hints already delivered can achieve the optimal route; a party
that reaches the finale missing Hint 3 cannot, regardless of their intent
or engagement level.

\subsection{The Advisory Period as Campaign Architecture}

The advisory-to-binding calibration --- Sessions 1--3 as advisory, Session 4
onward as binding --- is usually framed as a rubric-scoring decision (scores
in advisory sessions are not held to binding-mode thresholds). The evidence
supports reframing it as a campaign design decision. Sessions 1--3 are the
hint-planting window for Hints 1 and 2; they are the sessions in which PCs
encounter their spotlight material for the first time; they are the sessions
in which trust-accumulation systems are first loaded. Sessions 4--7 are the
sessions in which that material is used. A campaign that attempted binding-mode
quality in Session 1 would be structurally premature: the material for
binding-quality sessions depends on the planting work of Sessions 1--3.

The 4.5 mean sessions-to-first-binding-PASS validates this: in most campaigns,
the advisory period provides exactly the preparation the binding sessions require.
The exception (C4, first PASS at Session 6) is consistent with C4's extended
trust-building structure (7 watch-tokens, supply die attrition, HP carry-over):
C4 required more preparation sessions because its trust-mechanism was more
demanding.

\subsection{Route B as Designed Optimal: The C5 Case}

The professional-code archetype (C5, Thieves Guild) presents a result that
requires careful interpretation. The campaign achieved Route B, not D or E,
and Route B is listed in the route taxonomy as a ``personal-level resolution''
outcome rather than a ``full inversion'' outcome. Does this represent a spine
failure?

The answer is no, and the distinction matters. The C5 spine was designed with
Route B as the campaign's optimal achievable route. The guild-marker system
(trust tokens earned through professional acts, not personal-stakes acts)
was calibrated to a lower emotional register than the grief-themed or
covenant-themed campaigns. The central question (\textit{what does professional
loyalty cost?}) does not have a ``full inversion'' answer; it has a
``professional resolution'' answer, which is Route B. A Route D outcome for
C5 would have required personal-stakes material the archetype was explicitly
designed to exclude.

This indicates that the route-naming convention (A--E as difficulty-ordered)
should not be read as quality-ordered. Route E in a siege-resource campaign is
not inherently better than Route B in a professional-code campaign; they are
different kinds of answers to different kinds of questions. The spine's design
requirement is that the campaign's designated optimal route is achieved, not that
the letter designation is as high as possible.

\subsection{The Strangers Result and Archetype Independence}

The no-designed-stakes archetype (C7) achieving Route E is the corpus's most
consequential archetype-independence result. Earlier hypotheses about the spine's
scope held that the grief-hook structure and personal-stakes material in PCs'
character sheets were prerequisites for Route D/E completion. Campaign 7
falsifies this. The Reciprocal Acts log --- a system that records specific acts
of consideration between strangers without requiring grief-hooks or professional
codes --- produced Route E completion with an archetype that had neither.

The implication for the spine's generalizability is significant: the five
components do not require emotionally dense PC material to function. They require
\textit{some} form of trust-accumulation system calibrated to the archetype,
but the system's form is flexible. Grief-tokens and professional-markers and
reciprocal-acts are all instances of the same structural mechanism: specific
party actions earn specific evidence of engagement with the central question,
and that evidence gates the route conditions.

This suggests the spine is exportable to TRPG systems and settings with lower
personal-stakes norms than Dragonlance. Whether it functions equally well with
purely tactical parties (no social material at all) remains an open question;
it is the obvious next stress test.

\subsection{Limits of the Current Corpus}

Three limitations bound the current claims.

\textbf{AI simulation confound.} The 100\% completion rate is confounded by the
AI simulation context. No player attrition, no scheduling failures, no DM burnout.
Human-run campaigns using the spine would face all three. The spine's structural
design is a necessary but not sufficient condition for completion in human play;
this paper establishes sufficiency in the AI-simulated context only.

\textbf{Seven-session constraint.} All campaigns were designed for exactly 7 sessions.
This prevents generalization to shorter (5-session) or longer (10-session) campaigns.
The 7-session length was chosen as sufficient for full 4-hint delivery, 4-PC spotlight
distribution, and Route D/E conditions. Whether the spine's components scale to different
arc lengths is unexamined. In particular, a 5-session spine would need to compress the
hint delivery schedule (Hints 1 and 2 in Session 1; Hints 3 and 4 in Session 4), which
may reduce hint differentiation and convergence impact.

\textbf{Single designer.} All 7 campaign spines were designed by the same
researcher-designer. Cross-designer reliability --- whether other designers
implement the five components with equivalent outcomes --- is not established.
The component specifications are designed to be implementable from written
guidance, but this has not been tested with naive implementers.

\subsection{Future Work}

Three research directions follow directly from the corpus:

\textbf{Human-run validation.} The most important extension is running the spine
in human-DM tabletop play. The five components translate directly to human play;
the recovery-path mechanism is especially important in that context (human DMs
miss delivery windows more frequently than AI DMs following explicit condition
specifications). A human-run corpus of even 3--4 campaigns would substantially
strengthen the generalizability claim.

\textbf{5- and 10-session variants.} Testing the spine at arc lengths other than
7 would establish its flexibility bounds. The prediction is that 5-session arcs
will succeed with a compressed hint schedule but will produce lower Route D/E
rates; 10-session arcs will produce equivalent or higher Route D/E rates but risk
drift in sessions 6--9 (the window between Hint 4 delivery and the finale).

\textbf{Partial spine implementations.} The corpus contains no campaigns designed
with partial spine implementations (e.g., hints and routes but no deliverable
specification; deliverable and spotlights but no explicit end-conditions). A
controlled study varying which components are included would directly test the
component contribution analysis. The prediction --- that branching end-conditions
and hint schedule are necessary while spotlights and deliverable specification
are quality-enhancing rather than completion-necessary --- is testable.

The circularity concern raised by our operational definition of completion is
partially addressed by the three external indicators described in
Section~\ref{subsec:operationalizing-completion}: deliverable produced, hints
delivered, and Route D/E achieved. These indicators are observable without
rubric scoring and provide convergent-validity evidence for the Route D/E
classification. However, a fully independent measure of ``emotional completeness''
--- such as player-written post-campaign reflections documenting their
experience of the finale, or narrative coherence scoring by readers unfamiliar
with the rubric who assess whether the campaign's central question was answered ---
remains future work. Such independent measures would allow the rubric-based gate
score and the external-indicator composite to be validated against a criterion
that shares no measurement ancestry with either.
